\section{Related Work}

\subsection{XAttention}

\textbf{XAttention} (Xu et al., 2025) proposes \textbf{training-free} block scoring via 
antidiagonal sub-sampling. Blocks are ranked by the mean magnitude of their 
antidiagonal entries. We adopt the sub-sampled block scoring approach.

\subsection{SeerAttention}

\textbf{SeerAttention} (Gao et al., 2025) introduces a \textbf{learnable} gating network 
for block-level sparsity. It predicts per-head sparsity budgets using a 
low-rank query-key pooling network. We adopt a similar gating structure.

\subsection{DuSA}

\textbf{DuSA} (Cao et al., 2025) uses a \textbf{dual-stage} approach: dense intra-block 
attention (all tokens within a block) + sparse inter-block attention (top-$k$ 
blocks). We adopt this two-stage structure.

\subsection{AnchorAttention}

\textbf{AnchorAttention} (Zhu et al., 2025) applies \textbf{stripe band masking} at the 
token level. Tokens within a stripe group attend to adjacent groups. We adopt 
a simplified stripe mask.

\subsection{Spectraformer}

\textbf{Spectraformer} (Nguyen et al., 2025) uses RFF for global linear attention. 
We apply it periodically (every \texttt{rectify\_interval} layers) as a global 
rectification.

\subsection{Softpick}

\textbf{Softpick} (Zuhri et al., 2025) is a rectified softmax that produces 
natural sparsity (~47\% activation) without sink tokens. We use it in place 
of standard softmax.

\subsection{Position of Our Work}

\begin{table}[H]
\centering
\small
\begin{tabular}{llll}
\toprule
\textbf{Method} & \textbf{Learnable?} & \textbf{Training-free?} & \textbf{Complexity} \\
\midrule
XAttention & No & Yes & $\mathcal{O}(N B_s + N k B_s)$ \\
SeerAttention & Yes & No & $\mathcal{O}(N B_s + N k B_s)$ \\
DuSA & Partial & No & $\mathcal{O}(N B_s + N k B_s)$ \\
AnchorAttention & No & Yes & $\mathcal{O}(N S_w)$ \\
Spectraformer & No & Yes & $\mathcal{O}(N F)$ \\
\textbf{NEXUS} & \textbf{Yes} & \textbf{Partial} & \textbf{$\mathcal{O}(N B_s + N k B_s)$} \\
\bottomrule
\end{tabular}
\end{table}

\textbf{Key distinction.} NEXUS is the first to combine all six in a single 
module with modular enable/disable flags.
\begin{lstlisting}

---
\end{lstlisting}